\begin{trapbox}
	\begin{description}[style=nextline, leftmargin=2.8em, labelsep=0.5em, itemsep=0.35em]
		\item[\textbf{\textcolor{CTrap}{易错点一（公式类型混淆）}}] 将 $a^2-b^2$ 误认为 $(a-b)^2$，或直接展开成 $a^2-2ab+b^2$。
		\item[\textbf{\textcolor{CTrap}{正确理解（区分公式条件）：}}] $a^2-b^2$ 是平方差，需分解为 $(a+b)(a-b)$；$(a-b)^2$ 是完全平方，展开为 $a^2-2ab+b^2$，两者结构完全不同。
		\item[\textbf{\textcolor{CTrap}{错因分析（概念混淆）：}}] 只看到“平方”和“减号”就套用完全平方，忽略平方差公式要求“两个平方相减”且不含交叉项。
	\end{description}

	\medskip
	\begin{description}[style=nextline, leftmargin=2.8em, labelsep=0.5em, itemsep=0.35em]
		\item[\textbf{\textcolor{CTrap}{易错点二（完全平方缺项）}}] 展开 $(a+b)^2$ 时写成 $a^2+b^2$，缺少 $2ab$。
		\item[\textbf{\textcolor{CTrap}{正确理解（公式完整记忆）：}}] $(a+b)^2=a^2+2ab+b^2$，完全平方展开必须有三项，且中间项系数2。
		\item[\textbf{\textcolor{CTrap}{错因分析（记忆缺失）：}}] 错误认为平方运算直接分配到各项，忽略交叉项。
	\end{description}

	\medskip
	\begin{description}[style=nextline, leftmargin=2.8em, labelsep=0.5em, itemsep=0.35em]
		\item[\textbf{\textcolor{CTrap}{易错点三（立方和差符号）}}] 如将 $a^3+b^3$ 分解为 $(a+b)(a^2+ab+b^2)$，或 $a^3-b^3$ 分解为 $(a-b)(a^2-ab+b^2)$，二次项符号错误。
		\item[\textbf{\textcolor{CTrap}{正确理解（符号规律）：}}] $a^3+b^3=(a+b)(a^2-ab+b^2)$（中间减）；$a^3-b^3=(a-b)(a^2+ab+b^2)$（中间加）。
		\item[\textbf{\textcolor{CTrap}{错因分析（符号混淆）：}}] 未通过展开验证，仅凭记忆导致符号记反。
	\end{description}
\end{trapbox}
